\clearpage\begingroup
\SourceLeaf{151}\HistoricalHead{146}{L}{PHYSIK.\quad BRIEFWECHSEL.}
\begin{center}
[2.]\par\medskip
\textbf{\textsc{Gauss} an \textsc{Schumacher}.\quad Göttingen, 25. Junius 1810.}\par\medskip
\rule{.75\linewidth}{.4pt}\par
[Briefwechsel zwischen \textsc{Gauss} und \textsc{Schumacher} I, S. 45 f.]\par
\rule{.75\linewidth}{.4pt}
\end{center}
\noindent . . . . . . . . . . . . . . . . . . . . .\par
\noindent\hspace*{\parindent}\textsc{Repsolds} Wunsch, dass zwei von den vier Flächen einerlei Halbmesser bekommen, lässt sich nicht erfüllen, wenn die Bedingung Statt finden soll, dass die dritte Fläche einen kleinern Halbmesser habe, als die zweite. Sie können sich davon leicht selbst überzeugen, wenn Sie folgende 4 Formeln näher prüfen, wodurch nach \textsc{Repsolds} Angaben über die Brechungs- und Zerstreuungsverhältnisse die 4 Halbmesser $f,g,f',g'$ dargestellt werden, wenn die Brennweite 96 Zoll werden und die Abweichung, wegen Gestalt der Gläser und wegen der Farben gehoben werden sollen. Diese Form ist mir eigenthümlich[*)] und wie ich glaube, die zierlichste, die man finden kann:
\begingroup\renewcommand{\thefootnote}{}\footnotetext{[*) Vergl. \textit{Achromatische Doppelobjective, ohne Rücksicht auf Dicke und Abstand}, S. 135 dieses Bandes.]}\endgroup
\[
\begin{aligned}
\frac1f&=+0{,}0556255-0{,}0276427\operatorname{tang}\varphi\\[5pt]
\frac1g&=+0{,}0083843+0{,}0276427\operatorname{tang}\varphi\\[5pt]
\frac1{f'}&=+0{,}0129083-0{,}0339803\operatorname{secans}\varphi\\[5pt]
\frac1{g'}&=-0{,}0498917+0{,}0339803\operatorname{secans}\varphi.
\end{aligned}
\]

Der Winkel $\varphi$ kann nach Gefallen angenommen werden, negative Halbmesser zeigen \textls[120]{hohle} Flächen an.

Bei diesen Formeln ist die Dicke der Gläser nicht in Betrachtung gezogen. Mit Rücksicht auf diese werden einige Abänderungen nöthig; ich habe ein System von Werthen auf das schärfste berechnet, wodurch die äussern Farben bei den sehr nahe an der Axe [verlaufenden Strahlen] und zugleich die mittlern [Strahlen], die in einem Abstande von 2 Zoll von der Axe [verlaufen,] genau in einen Punkt zusammengebracht werden, und mir dabei die Bedingung vorgeschrieben, dass die 2\textsuperscript{te} und die 3\textsuperscript{te} Fläche nicht viel verschieden seyn sollen; je ungleicher man sie nimmt, desto kleiner wird der Halbmesser der ersten Fläche, und man wünscht kleine Halbmesser so viel [wie] möglich zu vermeiden.
\clearpage\SourceLeaf{152}\HistoricalHead{147}{R}{GAUSS, SCHUMACHER.}
\begin{center}
Maassen für ein Doppelobjectiv von 96 Zoll Brennweite,\\
5 Zoll Öffnung:
\end{center}
\begingroup\small
\[
\begin{array}{r@{\;}c@{\quad}c}
&&\text{Dicke in der Axe}\\[-2pt]
\begin{array}{r}\text{Halbmesser der Flächen des Convexglases}\\\vphantom{42}\end{array}
&\left.\begin{array}{r@{\,}l}26^{\mathrm{Zoll}}&{,}202\\42&{,}972\end{array}\right\}&0{,}21\\[7pt]
\begin{array}{r}\text{des Concavglases}\\\vphantom{100}\end{array}
&\left.\begin{array}{r@{\,}l}39&{,}985\\100&{,}845\end{array}\right\}&0{,}11
\end{array}
\]\endgroup

Ich bin überzeugt, dass, wenn die Brechungs- und Zerstreuungsverhältnisse genau so sind, wie \textsc{Repsold} sie gefunden hat, und die Flächen genau kugelförmig werden, dieses Glas eine sehr gute Wirkung thun muss. . . . . . . . .
\bigskip\begin{center}
[3.]\par\medskip
\textbf{\textsc{Schumacher} an \textsc{Gauss}[*)].}\par\medskip
\rule{.75\linewidth}{.4pt}\par
[Briefwechsel zwischen \textsc{Gauss} und \textsc{Schumacher} I, S. 52.]\par
\rule{.75\linewidth}{.4pt}
\end{center}
\begingroup\renewcommand{\thefootnote}{}\footnotetext{[*) Das Datum fehlt; dasselbe muß nach dem 5. August und vor dem 6. September 1810 liegen, wahrscheinlich wenige Tage vor dem letzteren.]}\endgroup
\noindent . . . . . . . . . . . . . . . . . . . . .\par
\noindent\hspace*{\parindent}\{\textsc{Repsold} hat jetzt nach vieler Mühe zwey Objective nach Ihrer Formel gefertigt. Das eine ist schlecht, das andere besser. Ganz lässt sich noch nicht darüber entscheiden, da es noch nicht sicher genug gefasst ist. Der Hauptgrund des Mislingens liegt wohl in der Dünne der Gläser, die so beym Poliren durchbiegen und ihre Gestalt verlieren. Die Brennweite trifft bis auf $0{,}005$ eines Fusses mit der berechneten überein. Er wird Ihnen noch das Minimum der Dicke der Gläser melden, über das sie nicht hinausgehen dürfen, ohne durchzubiegen, und dann um Modification Ihrer Formeln danach bitten.
\begin{flushright}d. 6\textsuperscript{ten} Sept. [1810].\end{flushright}
\noindent . . . . . . . . . . . . . . . . . . . . .\par
\noindent\hspace*{\parindent}\textsc{Repsold} hat jetzt die Fassung fertig, mit $\frac12$ Zoll Blendung thut es jetzt eben so viel, als das 6füssige, das Sie auf der hiesigen Sternwarte gesehen haben, aber auch nicht mehr, wie es doch wohl seyn sollte. Es ist immer ein vortreffliches Fernrohr, aber würde gewiss, wenn die Gläser nicht
\par\HistoricalFoot{}{19\textsuperscript{*}}
\clearpage\SourceLeaf{153}\HistoricalHead{148}{L}{PHYSIK.\quad BRIEFWECHSEL.}
\noindent durchgebogen wären, noch mehr thun. \textsc{Repsold} wünscht für die Decimaleintheilung des Fusses und der Öffnung die er machen kann, $=0^{\mathrm F}{,}38$, die Formeln so eingerichtet, dass das Convexglas in der Axe
\begin{center}$0^{\mathrm F}{,}025$ dick seyn kann,\end{center}
\noindent und das Concavglas ebenda
\begin{center}$0^{\mathrm F}{,}018$ seyn kann;\end{center}
\noindent statt dass z. B. bey diesem letzten Sie bey der Berechnung nur $0^{\mathrm F}{,}009$ angenommen haben.

Das jetzt vollendete Fernrohr hat in der Axe
\begin{center}
bey dem convexen $0^{\mathrm F}{,}021$,\\
bey dem concaven $0^{\mathrm F}{,}012$,
\end{center}
\noindent die Brennweite $=8^{\mathrm F}{,}06$. Ein wenig färbt es noch bey Fixsternen. . . . . . .\}
\bigskip\begin{center}
[4.]\par\medskip
\textbf{\textsc{Gauss} an \textsc{Schumacher}.\quad Göttingen, 6. Oktober 1810.}\par\medskip
\rule{.75\linewidth}{.4pt}\par
[Briefwechsel zwischen \textsc{Gauss} und \textsc{Schumacher} I, S. 54.]\par
\rule{.75\linewidth}{.4pt}
\end{center}
\noindent . . . . . . . . . . . . . . . . . . . . .\par
\noindent\hspace*{\parindent}Dass das Objectiv, welches \textsc{Repsold} geschliffen, nicht ganz nach seiner Erwartung ausgefallen ist, thut mir leid. Ich will die Rechnung noch einmal für eine etwas grössere Dicke der Gläser wiederholen, wozu ich bisher noch nicht habe kommen können, da ich die auf meine erste Rechnung sich beziehenden Papiere verlegt und bisher nicht Zeit gehabt habe, mich ganz von neuem wieder hineinzusetzen. Ich glaube aber, dass eine etwas vergrösserte Dicke die Dimensionen nur wenig verändern wird. . . . . . . . . . . . . . . .
\bigskip\begin{center}
[5.]\par\medskip
\textbf{\textsc{Schumacher} an \textsc{Gauss}.\quad Altona, 10. November 1810.}\par\medskip
\rule{.75\linewidth}{.4pt}\par
[Briefwechsel zwischen \textsc{Gauss} und \textsc{Schumacher} I, S. 58.]\par
\rule{.75\linewidth}{.4pt}
\end{center}
\noindent\hspace*{\parindent}\{\textsc{Repsold} hat jetzt nach Ihren Formeln, werthester Freund, ein zweytes Objectiv vollendet, das ohnerachtet das Glas Streifen hat, vortreffliche
\clearpage\SourceLeaf{154}\HistoricalHead{149}{R}{GAUSS, SCHUMACHER, OLBERS.}
\noindent Wirkung thut. Es sitzt schon im Passageninstrumente und zeigt bey hellem Tage Mizar als Doppelstern. Beyde Sterne sind keine $20''$ von einander entfernt. Alle Sterne der vierten Grösse kann man jetzt bey Tage beobachten. Er ist voll Dank gegen Sie, muss aber bald, da er französisches Flintglas von \textsc{Cruynes} und \textsc{Lauçon} erhält, unterthänigst bey Ihnen einkommen und um neue Maasse, nach dem neuen Brechungs- und Zerstreuungsverhältniss, bitten? Wollen Sie nicht Ihre Formeln bekannt machen? . . . . . . .\}
\bigskip\begin{center}
[6.]\par\medskip
\textbf{\textsc{Gauss} an \textsc{Schumacher}.\quad Göttingen, 6. Januar 1811.}\par\medskip
\rule{.75\linewidth}{.4pt}\par
[Briefwechsel zwischen \textsc{Gauss} und \textsc{Schumacher} I, S. 61.]\par
\rule{.75\linewidth}{.4pt}
\end{center}
\noindent . . . . . . . . . . . . . . . . . . . . .\par
\noindent\hspace*{\parindent}Dass \textsc{Repsolds} Glas so gut ausgefallen ist, freuet mich sehr. Wenn er französisches Flintglas erhalten und dessen Zerstreuungskraft durch Versuche geprüft hat, so will ich gern die Krümmungen der Linsen daraus berechnen, allein ich muss dann bitten, mir zugleich die Dimensionen des nach meinen Zahlen geschliffenen Glases wieder mitzutheilen, da ich sie alles Suchens ungeachtet nicht habe wieder auffinden können, und sonst wenn ich ganz von vorn anfangen müsste, eine viel längere Arbeit haben würde. . . . . . . . . .
\clearpage\endgroup
